\section{部分積分の反復}\label{節:部分積分の反復}

前節までは,一度の部分積分で既知の積分に到達した.では,整式の次数が高い場合はどうなるだろうか.一般的な形は,\,$P(x)$ を $n$ 次の整式,\,$a\neq0$ として\footnote{$a=0$ なら $\int P(x)\dd x$ であり,部分積分は要らない.}
\[
 \int P(x)\e^{ax}\dd x,\qquad
 \int P(x)\cos ax\dd x,\qquad
 \int P(x)\sin ax\dd x
\]
で与えられる.整式を微分する側に選ぶと,一度の部分積分で整式の次数が1つ下がる.これを繰り返して整式を定数にし,\,$\int\e^{ax}\dd x$ などの既知の積分に帰着させることを目標とする.

指数関数の場合を書き下してみよう.\,$u=P(x)$,\,$v'=\e^{ax}$ と選ぶと,\,$v=\frac1a\e^{ax}$ だから
\[
 \int P(x)\e^{ax}\dd x
 =\frac1aP(x)\e^{ax}-\frac1a\int P'(x)\e^{ax}\dd x
\]
となる.右辺の積分は,整式の次数が1つ下がっただけで,もとと同じ形である.そこで同じ部分積分をもう一度施してやれば,
\[
 \int P(x)\e^{ax}\dd x
 =\frac1aP(x)\e^{ax}-\frac1{a^2}P'(x)\e^{ax}
  +\frac1{a^2}\int P''(x)\e^{ax}\dd x
\]
となる.ここで,\,$P$ を $k$ 回微分した関数を\textbf{第$k$次導関数}といい,\,$P^{(k)}$ と書く.\,$P^{(1)}=P'$,\,$P^{(2)}=P''$ である.\,$n$ 次の整式は微分するたびに次数が1つ下がるので,\,$P,P',P'',\ldots$ と微分を繰り返すと,第$n$次導関数 $P^{(n)}$ は定数になる.したがって $n$ 回目の部分積分の後に残る積分は $\int\e^{ax}\dd x$ の定数倍であり,ここで計算が終わるのだ.まとめると
\[
 \int P(x)\e^{ax}\dd x
 =\e^{ax}\left\{\frac{P(x)}{a}-\frac{P'(x)}{a^2}+\frac{P''(x)}{a^3}-\cdots
  +(-1)^n\frac{P^{(n)}(x)}{a^{n+1}}\right\}+C
\]
を得る.

三角関数の場合も同じ手順で進む.異なるのは,\,$\cos ax$ と $\sin ax$ を積分するたびに $\frac1a$ が掛かるだけでなく,正弦と余弦が入れ替わり,符号も変わることである.次数が下がることと,係数・符号が変わることを別々に確認すればよい.

\begin{bookpoint}
\textbf{POINT}\par
整式を微分する側,指数関数・三角関数を積分する側に選び,整式が定数になるまで部分積分を繰り返す.
\end{bookpoint}

\begin{example}{部分積分の反復}
$\int x^2\cos x\dd x$ を求めよ.
\end{example}

\paragraphbf{方針}
$x^2$ を微分し,\,$\cos x$ を積分する.残る $\int x\sin x\dd x$ でもう一度,整式を微分する部分積分を行う.

\paragraphbf{解答}
$x^2\sin x+2x\cos x-2\sin x+C$

\paragraphbf{解法}
$u=x^2$,\,$v'=\cos x$ と選ぶ.\,$u'=2x$,\,$v=\sin x$ なので,
\[
 \int x^2\cos x\dd x
 =x^2\sin x-2\int x\sin x\dd x
\]
となる.残った積分では,\,$x$ を微分し,\,$\sin x$ を積分する.\,$\sin x$ の原始関数は $-\cos x$ だから,
\[
\begin{aligned}
 \int x\sin x\dd x
 &=x(-\cos x)-\int1\cdot(-\cos x)\dd x\\
 &=-x\cos x+\int\cos x\dd x\\
 &=-x\cos x+\sin x+C
\end{aligned}
\]
である.これをもとの式へ戻して,任意定数をまとめると
\[
\begin{aligned}
 \int x^2\cos x\dd x
 &=x^2\sin x-2(-x\cos x+\sin x)+C\\
 &=x^2\sin x+2x\cos x-2\sin x+C
\end{aligned}
\]
を得る.2回目の積分で現れた負号だけでなく,その式全体に最初の $-2$ が掛かることに注意する.途中の積分を独立に計算したときは,戻す先の係数を残しておくと,符号の取り違えを防げる.

\paragraphbf{答えの微分による確認}
答えを微分すると,途中で現れた項が対になって消える.
\[
\begin{aligned}
 &\frac{\dd}{\dd x}
  \{x^2\sin x+2x\cos x-2\sin x\}\\
 &\qquad=2x\sin x+x^2\cos x
       +2\cos x-2x\sin x-2\cos x\\
 &\qquad=x^2\cos x
\end{aligned}
\]
$x^2\sin x$ の微分で生まれた余分な項 $2x\sin x$ を,\,$2x\cos x$ の微分の $-2x\sin x$ が打ち消す.その際に生まれた $2\cos x$ を,さらに $-2\sin x$ の微分が打ち消す.部分積分で補った項は,このように積の微分で生まれた余分な項を順に打ち消しているのである.

\paragraphbf{積分の漸化式との関係}
次数を文字のまま残しておくこともできる.\,$I_n=\int x^n\e^{ax}\dd x$ と置き,\,$x^n$ を微分する部分積分を1回だけ行えば,
\[
 I_n=\frac1ax^n\e^{ax}-\frac naI_{n-1}
\]
となる.これは $I_n$ を $I_{n-1}$ で表す漸化式である.このように,部分積分で次数を1つ下げる操作は,積分を項とする数列の漸化式として書き下せる.\chapterref{章:積分と数列}では,この形の式を使って定積分の値を求める.
